% Generated draft: requires mathematical and source review.
\documentclass[11pt]{article}
\usepackage[margin=25mm]{geometry}
\usepackage{fontspec}
\setmainfont{Latin Modern Roman}
\usepackage{amsmath,amssymb,mathtools,unicode-math}
\setmathfont{Latin Modern Math}
\usepackage{enumitem,needspace,xurl,hyperref}
\setlength{\parindent}{0pt}
\setlength{\parskip}{5pt}
\setlength{\emergencystretch}{3em}
\newcommand{\N}{\mathbb{N}}
\newcommand{\Z}{\mathbb{Z}}
\newcommand{\Q}{\mathbb{Q}}
\newcommand{\R}{\mathbb{R}}
\newcommand{\C}{\mathbb{C}}
\newcommand{\F}{\mathbb{F}}
\newcommand{\A}{\mathbb{A}}
\newcommand{\PP}{\mathbb{P}}
\newcommand{\E}{\mathbb{E}}
\newcommand{\eps}{\varepsilon}
\newcommand{\dd}{\,\mathrm{d}}
\begin{document}
\section*{UnsolvedMath ID: 8700055}
\textbf{Source problem code:} catalogue entry
\subsection{Definitions and mathematical statement}
Let $m\geq1$ be an integer and $\theta\in\C^m$. The Laurent--Roy conjecture in the source asserts that there exists $c=c(\theta,m)>0$ such that, for every integer $k$ with $0\leq k\leq m$, infinitely many integers $T\geq1$ admit a $\Q$-defined algebraic set $Z\subset\C^m$ of dimension exactly $k$ and a point $\alpha\in Z$ satisfying
\[
t(Z)\leq T^{m-k},\qquad
|\theta-\alpha|\leq\exp\{-cT^{m+1}\}.
\]
The constant is chosen before $k$ and is independent of $T,Z,\alpha$. The infinitely many $T$ may depend on $k$; the source does not require one common infinite set of $T$ working for every $k$, nor every sufficiently large $T$.

An algebraic set defined over $\Q$ means a common zero locus of polynomials in $\Q[x_1,\ldots,x_m]$, with its usual algebraic dimension. No irreducibility or equidimensionality assumption is supplied. The point $\alpha$ is an arbitrary complex point on $Z$, not stated to be algebraic over $\Q$ in positive dimension. The norm $|\cdot|$ is unspecified; a fixed Euclidean norm is a conventional interpretation requiring source confirmation.

The crucial size $t(Z)$ is not defined in the record. It must be recovered, including its degree/height and affine/projective normalizations if applicable; none is guessed. This affects the endpoint $k=m$, where the required bound is $t(Z)\leq1$. Both endpoints $k=0,m$ are retained. The summary's mention of hypersurfaces does not reduce the source's full range of dimensions to $k=m-1$.
\subsection{Short English statement}
Approximate a fixed complex vector exponentially closely by a point on a rationally defined algebraic set of each prescribed dimension, with the exact size bound $t(Z)\leq T^{m-k}$. The size invariant is missing, but the uniform constant and infinitely-many-parameter quantifiers are explicit.
\subsection{Sources}
\begin{itemize}
\item Local UnsolvedMath registry, permanent ID 8700055, exact source code catalogue entry, ``Conjecture 4.13---Laurent--Roy''; registry attribution Michel Waldschmidt, no proposed year supplied. Registry status open, classification OPEN-TRIAGE.
\item Michel Waldschmidt, \emph{Open Diophantine Problems} (2004), Conjecture 4.13, \url{https://arxiv.org/abs/math/0312440}. Supplied source-TeX extraction accessed July 29, 2026.
\item The supplied summary attributes the conjecture to Laurent--Roy and mentions related weaker/transposed approximation bounds by Roy and Laurent--Roy without precise references. It reports the exact sharp statement open in its accessed literature; this is not an independent current-status check.
\item Provenance flags current primary-literature status and redistribution terms NEEDS\_REVIEW; the historical survey presents numbered open problems and conjectures.
\end{itemize}
\end{document}
